\chapter{Masa Lalu}
\penulis{M Raditya Arka Fakhlevi}
\awalcerpen{K}{enalin aku raka}, dan di sini aku akan menceritakan masa lalu aku. \\


Jadi $\dots$

Mau kah kamu mendengarkan cerita ku?


Kalau ya, ayo kita mulai


Ini adalah kisah saat aku TK


Saat TK aku termasuk orang yang culun sering nangis, penakut, dan cuman bisa mengandalkan orang lain, dan kadang selalu di hina bahkan aku pernah di TK ada kejadian dimana aku di hina oleh 5 orang termasuk orang yang aku sukai saat itu. Kenapa aku tidak melawan? itu lucu, karena salah satu dari anak tersebut, salah satu orang berada dan dia memiliki badan seperti anak SD. Yang aku lakukan hanyalah menangis di belakang guru, hanya itu. Itu membuatku selalu melihat ke bawah, tidak berani melihat wajah orang yang sedang berbicara dengan ku sampai itu terus berjalan sampai SD.


Setelah aku berumur 7 tahun aku pun pindah rumah, tanpa berpamitan. Dengan semua temanku.


Kamu masih membacanya? Kalau ya akan aku lanjut


Ini adalah cerita ku di SD


Saat itu di kelas 2 SD


Hari pertama ku sekolah, aku sangat bersemangat untuk mendapatkan teman tapi aku hanya duduk diam dan menangis karena aku kurang pandai bersosialisasi dan saat anak-anak lain sudah bisa mengobrol dan berkenalan satu sama lain, yang aku lakukan hanyalah menangis karena tidak dapat berkenalan, padahal saat di TK dan kelas 1 SD aku termasuk orang yang mudah bersosialisasi tapi kenapa di kelas 2, di sekolah yang baru, aku tidak bisa mudah bersosialisasi lagi sampai ibuku sendiri datang untuk menenangkan diriku dan membantuku mendapatkan teman.


Memalukan? Tentu saja, anak yang berusia 8 tahun nangis di tenangin ibunya di liat banyak orang, mana mungkin bisa dilupain.


Dan terkadang saat ibuku menjemput dia selalu melihat nilai di buku ku dan melihat tulisan ku, dan $\dots$


Ibuku selalu bilang : ``Kenapa nilai-nya merah. Kan sudah di suruh belajar malah main da ini mah anak-nya pinter banget.'' Di depan ibu anak-anak lain, mungkin di luar aku tersenyum tapi di dalam aku merasa bersalah dan aku harus kuat karena di sindir sama ibu-ku sendiri. Tapi aku terima-terima aja jika itu faktanya dan sampai sekarang pun masih seperti itu.


Hahahhaha.


Ok mari lanjut, di kelas 4, kenapa langsung ke kelas 4 karna tidak ada kejadian menarik di kelas 3.


Ok lanjut di kelas 4


Saat itu ada anak baru katanya di pindahan dari sekolah lain dan juga dia baru pindah rumah juga, persis seperti diriku, aku pun ingin mulai bisa bersosialisasi kembali agar bisa mudah berteman kembali. saat itu aku tidak sempat berkenalan atau mengobrol dengannya. Tapi saat keesokannya saat ingin berkenalan pada jam olahraga dia malah menurunkan celana ku.


Lucu? Tentu saja hahahh untuk orang yang menurunkan celananya dan melihatnya. Tapi tidak untuk korbannya.


Oke saat itu aku ikut tertawa tapi tidak dengan diriku di dalam apa yang aku rasakan, tentu saja malu, tapi aku berusaha tidak menangis. Baik kita lanjut.


Beberapa minggu kemudian saat itu sedang praktek di pelajaran olahraga saat aku sedang berbaris dan menunggu giliran 2 orang di belakang ku bagas dan fadhil mereka malah menyentil telinga ku secara bergantian,


Mungkin bagi mereka itu cuman candaan tapi untuk ku itu bukan candaan.


Bahkan daun telinga ku sampai merah aku pun bilang buat berhenti tapi bukannya berhenti, saat aku melihat kedepan mereka malah menurunkan celana ku. Aku berusaha menghiraukan nya tapi mereka terus melakukan nya aku pun sudah tidak kuat, akupun menangis saat aku di bawa ke kelas sekilas aku melihat guru olahraga ku dia hanya diam meng absen murid yang sudah praktek. Saat itu aku mulai membenci guru itu bahkan sampai sekarang.


Kenapa aku masih membenci guru itu? Karena


Katanya guru itu orang tua ke 2 bagi anak tapi kenapa saat melihat anak nya menangis dia tampak tak peduli, mungkin kalian cuman bilang ``ah cuman gitu doang masa di benci segitunya'', ok ga masalah, karna itu pendapat ku karena setiap orang memiliki pendapat dan perspektif yang berbeda beda. Ok mungkin kita sudahi saja cerita di kelas 4.


Mari kita lanjutkan cerita ini ke kelas 6


loh kok kelas 6 sih?


Ya mau gimana lagi pas kelas 5 kan ada pandemi covid-19. Di kasih tau libur 2 minggu eh ternyata 2 tahun. Ya udah mari lanjut cerita nya


Cerita kelas 6


Sebenarnya di sini ga ada yang bisa banyak di ceritain tapi buat ku ini salah satu yang penting


Setelah pandemi mereda aku dan semua murid di SD kembali bersekolah dan apa yang berubah? Ya sikap semua orang termasuk teman ku, panggil aja Qila pas di kelas 4 dia itu orang yang sopan dan ga pernah ngomong kasar tapi setelah masuk kembali dia berubah benget, jadi sering ngomong kasar, atau menghina, melihat perubahan ini aku terkejut berpisah 2 tahun banyak perubahan sikap.


Jadi, apa yang harus aku lakukan?


Mengikuti arus = \textit{ikut ikutan berkata kasar dan jarang di temani dia.} Atau Melawan arus = \textit{tidak berkata kasar tapi harus berusaha lebih keras agar tetap bisa berteman.}


Ya bimbang itu yang aku rasakan, tapi aku memilih opsi ke 2 walau harus lebih keras agar bisa dianggap tapi setidaknya aku masih bisa berbicara sopan di hadapan orang lain


Walaupun aku memilih opsi ke 2 tapi setidaknya aku masih bisa berteman dengan Qila sampai lulus sd dan sampai sekarang kita masih saling berhubungan walau sudah jarang bertemu.


Hmmmm$\dots$


Mungkin sudah cukup aku bercerita, aku bisa aja lanjut bercerita tapi aku memiliki alasan berhenti bercerita. 
Jadii apa yang aku dapat kan selama di SD


Kenangan buruk? Mungkin


Pelajaran? Tentu saja


Meski menurut kamu di cerita itu aku terlihat seperti lebay atau cengeng menurut ku setiap anak memiliki mental yang berbeda, kenapa anak laki laki tidak boleh menangis?


Apakah ada peraturan dimana seorang laki laki tidak boleh nangis?


Candaan seperti itu bukan membuat mental anak kuat tapi bisa jadi membuat kejadian itu menjadi trauma yang dalam dan tidak bisa di lupakan, mungkin saat bertemu, kita akan seperti biasa tapi belum tentu bisa melupakan masa lalu yang buruk. Dan untuk orang dewasa jangan anggap sepele hal seperti itu mungkin itu seperti kejahilan kecil atau main main tapi kadang kejahilan seperti itu bisa berdampak buruk.


Dan apa yang aku alami membuat ku takut melihat wajah orang sulit bersosialisasi dan juga tidak mau bercerita kepada orang lain


Dan alasan aku memakai kisah nyata karna aku sudah bisa berdamai dengan diri sendiri walaupun kadang masih sulit untuk percaya kepada diri sendiri.
